\scanpage{216}
\begin{center}\uuline{POINTS DE GORENSTEIN D'UN MORPHISME}

par Manuela STOIA\end{center}

R\deleted{\unsure{e|a}}appelons qu'un anneau noethérien est de Gorenstein s'il est de dimension injective finie ( \handwritten{3} ). On va utiliser
le critère suivant: ``\uline{Un anneau local noetherien $\underline A$,quotient d'un anneau local noethérien régulier $\underline B$, est de Gorenstein
si et seulement si $\operatorname{Ext}^{\underline q}_{\underline B}(\underline A,\underline B)=0$ pour tous les entiers $\underline q$ à
l'exception d'un $\underline q_{\underline{\unsure{o|0}}}$ et $\operatorname{Ext}^{\handwritten{\underline q_{\underline o}}}_{\underline B}(\underline A,\underline B)\simeq\underline A$}''\handwritten{[(3),prop\unsure{.|:}1.3,ch 1,§ 1.4]}.

Soit $\underline k$ un corps algébriquement clos.On va considérer des variétés algébriques $\underline X$ sur $\underline k$ au sens de ( \handwritten{4} ), mais pouvant avoir des éléments nilpotents.
On dit que $\underline X$ est\unsure{:|.} de Gorenstein
au point $\underline x$ (réspect.de Gorenstein) si l'anneau
local $\handwritten{\mathscr O_{\underline X,\underline x}}$ est de Gorenstein (réspect.si tous les anneaux $\handwritten{\mathscr O_{\underline X,\underline x}}$ sont de Gorenstein).

\uline{Théorème}: \uline{Si $\underline f:\underline X\handwritten{\longrightarrow}\underline Y$ est un morphism\handwritten{e} plat de variétés algébriques, alors l'ensemble} :
\[
\underline{\underline G_{\handwritten{\underline f}}(\handwritten{\underline{\unsure{X|x}}})=\handwritten{\bigl\{}\handwritten{\underline x\in}\underline X/\underline X_{\handwritten{\underline f(\underline x)}}\quad\text{soit de Gorenstein au point }\handwritten{\underline x\bigr\}}}
\]
\uline{est ouvert dans $\underline X$ (pour $\underline y\handwritten{\in}\underline Y$,on note comme d'habitude par $\underline X_{\underline y}$
la variété algébrique- fibre d\deleted{\unsure{a}}e $\underline f$\deleted{\unsure{,}} au-dess\deleted{o}us de $\underline y$)}.

\uline{Démonstration}. L'énoncé est local sur $\underline X$,donc on peut
supposer $\underline X$ affine et soit $\deleted{\unsure{i|f}}\handwritten{\underline j}:\underline X\handwritten{\lhook\joinrel\longrightarrow}\handwritten{\underline k}^{\underline n}$ une immersion fermée.
Notons par $\underline i$ l'immersion donné par la composition
\diagram{\begin{tikzcd}[ampersand replacement=\&,column sep=large]
\underline X\arrow[r,"{\handwritten{\text{\textcircled{$\Gamma$}}_{\underline f}}}"]\&\underline X\times\underline Y\arrow[r,"{\handwritten{\underline j\times\underline 1_{\underline Y}}}"]\&\handwritten{\underline k}^{\underline n}\times\underline Y
\end{tikzcd}}
et par $\handwritten{\text{\textcircled{$\pi$}}}:\handwritten{\underline k}^{\underline n}\times\underline Y\handwritten{\longrightarrow}\underline Y$ la projection.On a la rélation $\underline f=\handwritten{\text{\textcircled{$\pi$}}}\circ\underline i$. Pour la commodité,
% archive sheet number: 215
